% Auxiliary module calculation supplied to maps for the complement question.
Let $\widehat W\to W$ be the infinite cyclic cover in which $\gamma$
acts by $z$ and $u$ acts trivially, and put $R=\mathbb Z[z,z^{-1}]$.
Then its ordinary homology, as integral deck modules, is
\[
 H_0=H_1=R/(z-1),\qquad H_2=R\oplus(R/(z-1))^2,
 \qquad H_3=0,\qquad H_4=R.
\]
Indeed the same normalization descent as D.1.2 works with direct sums
of strata, since every ordinary singular chain has compact support.
The degree-zero row is the cellular chain complex of the infinite
cylinder which covers the dual torus. Its homology is $R/(z-1)$ in
degrees zero and one and zero in degree two. The degree-four row is
$R$ in normalization degree zero. The degree-two homology differential
from double curves to normalized components is the transpose, up to
the unimodular intersection marking and the possible inversion
$z\mapsto z^{-1}$, of
\[
 B(z)=\begin{pmatrix}-1&-1&-1&-1\\1&1&1&z\\-z&-z&-1&-z\end{pmatrix}.
\]
The explicit unimodular operations in
\texttt{code/manifolds\_cusp\_alexander.py} reduce this to
$\operatorname{diag}(1,z-1,z-1)$ with one extra source column.
Its transpose is injective over the integral domain $R$, and its
cokernel is $R\oplus(R/(z-1))^2$. There are no higher differentials
because the normalization direction has degrees only zero through two
and all odd homology of the strata vanishes. There is no extension in
total degree two: the other possible piece, degree-two homology of the
dual cylinder, is zero. This proves the integral module assertion.
As a consistency check, derived reduction modulo $z^d-1$ gives the
finite-cover ranks $H_2=\mathbb Z^{d+3}$ and $H_3=\mathbb Z^2$, including
the additional Tor contribution from $H_1$; these are exactly D.1.2.
